\begin{figure}[p]
\centering
\begin{tikzpicture}[x=1.7cm,y=1.7cm,font=\scriptsize]
\draw[gray,dashed] (0,0) rectangle (3,3);
\foreach \y in {0,1,2}{\foreach \x in {0,1,2}{
 \pgfmathtruncatemacro{qh}{3*\y+\x}
 \pgfmathtruncatemacro{qv}{9+3*\y+\x}
 \draw[teal!70!black,thick] (\x,\y)--(\x+1,\y);
 \draw[teal!70!black,thick] (\x,\y)--(\x,\y+1);
 \node[circle,draw=teal!70!black,fill=white,inner sep=2.5pt] at (\x+.5,\y) {};
 \node[below=5pt] at (\x+.5,\y) {$q_{\qh}$};
 \node[circle,draw=teal!70!black,fill=white,inner sep=2.5pt] at (\x,\y+.5) {};
 \node[left=5pt] at (\x,\y+.5) {$q_{\qv}$};
}}
\foreach \y in {0,1,2,3}{\foreach \x in {0,1,2,3}{
 \pgfmathtruncatemacro{v}{3*mod(\y,3)+mod(\x,3)}
 \fill (\x,\y) circle (2pt);
 \node[above right=2pt] at (\x,\y) {$v_{\v}$};
}}
\draw[<->,gray] (.4,-.55)--(2.6,-.55);
\node at (1.5,-.8) {Horizontal boundary identification};
\node[rotate=90] at (3.65,1.5) {Vertical boundary identification};
\end{tikzpicture}
\caption{A fundamental-domain chart of the actual $3\times3$ toric-code manifest. Boundary vertex labels repeat because opposite sides are identified. The eighteen white edge midpoints, not the thirty-six incidence darts, index the physical qubits. The selected first star is $\{q_0,q_2,q_9,q_{15}\}$.}
\end{figure}
